\SourceLeaf{125}
\fancyhead[L]{\small General Theory of Terrestrial Magnetism}
\renewcommand{\Article}[1]{\par\medskip\Needspace{4\baselineskip}\pdfbookmark[1]{Art. #1}{erdmagnetismus-art-#1}\centerline{\textbf{#1.}}\nobreak\smallskip}
\pdfbookmark[0]{General Theory of Terrestrial Magnetism}{erdmagnetismus-body}
\par\medskip
\centerline{GENERAL THEORY}
\centerline{OF}
\centerline{TERRESTRIAL MAGNETISM.}
\medskip
The tireless zeal with which people have lately sought to investigate the direction and strength of the Earth's magnetic force in every part of its surface is all the more heartening as the purely scientific interest behind it becomes clearer. For navigation, it is indeed important to know the line of declination as thoroughly as possible. Navigation requires nothing beyond that and has scarcely any interest in the rest. Science gladly serves material interests as well, but its scope is wider and it calls for equal effort on every aspect of its inquiry.

The results of magnetic observations are usually represented on maps of the Earth by three systems of lines, which have been called isogonic, isoclinic, and isodynamic lines. Their shapes and positions change considerably over time, so a single drawing shows the state of the phenomenon only at a particular moment. \textsc{Halley}'s declination map differs greatly from \textsc{Barlow}'s representation in 1833, and \textsc{Hansteen}'s inclination map for 1780 already differs greatly from the present position of the isoclinic lines. Attempts to map the intensity are too recent for similar changes \SourceLeaf{126}to be demonstrated in it yet, though such changes will undoubtedly occur in time. All these maps still have gaps to a greater or lesser extent, or are unreliable in places. Even if the inaccessibility of some parts of the Earth's surface keeps the maps from becoming quite complete, there is hope that they will come much closer to completeness in rapid strides.

Viewed from the higher standpoint of science, however, assembling as complete an account of the phenomena as observation allows is still short of the actual goal. It is much like what an astronomer has done after observing, for example, the apparent path of a comet on the celestial sphere. There are only building stones and no building until the intricate phenomena have been brought under a single principle. The astronomer's main work begins only after the celestial body has passed out of sight. Relying on the law of gravitation, the astronomer calculates the elements of its true orbit from the observations and can even predict its further course with certainty. In the same way, the physicist should undertake, at least as far as the different and in some respects less favorable circumstances permit, to investigate the fundamental forces that produce the phenomena of terrestrial magnetism, both how they act and what their magnitudes are. The physicist should bring the observations, as far as they extend, under these elements and thus be able to anticipate the phenomena, at least with some degree of reliable approximation, in regions that observation has been unable to reach. In any event, it is worth keeping this highest goal in sight and trying to open up the paths that lead to it, even though the great imperfection of the available data still allows only a very rough approximation to the goal itself.

I do not intend to discuss the earlier unsuccessful attempts whose authors thought they could guess the answer to the great riddle without any physical basis. A physical basis can be granted only to attempts that treat the Earth as a real magnet and base the calculation solely on the established way a magnet acts at a distance. But all attempts of this kind so far have one feature in common. Instead of first asking \textit{how} this great magnet must be constituted to account for the phenomena, with an open mind as to whether its constitution would prove \SourceLeaf{127}simple or highly complex, they took a particular simple constitution as their starting point and tested whether the phenomena were compatible with that hypothesis. This merely repeats what the history of astronomy and the natural sciences tells us about the beginnings of so much of our knowledge.

The simplest hypothesis of this kind assumes a single very small magnet at the Earth's center. More precisely, since hardly anyone seriously believed such a magnet actually existed, it assumes that magnetism is distributed through the Earth so that its total outward action is equivalent to that of an imaginary infinitely small magnet. This is much like gravitation toward a homogeneous sphere being equivalent to attraction by an equal mass concentrated at its center. On this assumption, the Earth's magnetic poles are the two points where the extended magnetic axis of that central magnet intersects the Earth's surface. At these points the magnetic needle stands vertical and the intensity is greatest. On the great circle midway between the two poles, the magnetic equator, the inclination is $=0$ and the intensity is half what it is at the poles. Between the magnetic equator and a pole, both inclination and intensity depend only on the distance from that equator, the magnetic latitude, with the tangent of the inclination equal to twice the tangent of this latitude. Finally, the horizontal needle everywhere points along a great circle drawn toward the northern magnetic pole. Nature agrees with all these necessary consequences of the hypothesis only roughly. In reality the line of zero inclination is a curve of double curvature rather than a great circle. Equal inclinations do not give equal intensities. The directions of the horizontal needle are far from all converging on a single point, and so on. Even a cursory look therefore suffices to show that the hypothesis must be rejected. Yet one of the propositions above is still used as an approximation to determine the position of the line of zero inclination from observations made some distance from it, where the inclinations are moderate.

Eighty years ago \textsc{Tobias Mayer} had already \SourceLeaf{128}started from a similar hypothesis, with the difference that he placed the infinitely small magnet about one seventh of the Earth's radius from the center. Presumably to avoid complicating the calculation further, he retained the otherwise wholly arbitrary restriction that the plane perpendicular to the magnet's axis should pass through the Earth's center. In this way he found quite good agreement between his calculation and the observed declinations and inclinations, admittedly at only a very small number of places. A more extensive test would soon have shown, however, that this hypothesis could represent the full range of phenomena of these two elements little better than the first one. As is well known, there were as yet no intensity measurements at all.

\expandafter\def\csname HistoricalNoteText14\endcsname{For declination there is even one case of a difference exceeding 29 degrees. It is only fair, however, to judge the error of the calculation by the actual difference between the calculated and observed directions as a whole. The number of degrees of declination by itself does not give that difference, which is $11\frac{1}{2}$ degrees at the place in question.}
\textsc{Hansteen} went a step further by trying to fit the hypothesis of \textit{two} infinitely small magnets of unequal position and strength to the phenomena. The decisive test of whether a hypothesis is admissible is always to compare the results it yields with experience. \textsc{Hansteen} compared his hypothesis with observations at 48 different places. Intensity was also determined at only 12 of them, and all three elements are available at only 6. We still find discrepancies between calculation and observation, reaching nearly 13 degrees in inclination\HistoricalNote{14}.

If we find discrepancies of this size incompatible with what a satisfactory theory must provide, we cannot avoid concluding that the Earth's magnetic constitution cannot be represented by a concentration in one or a pair of individual infinitely small magnets. It remains possible that a \textit{greater} number of such imaginary magnets might eventually yield satisfactory agreement. Whether it would be advisable to solve the problem in that form is an entirely different question. Indeed, the calculations, already so \SourceLeaf{129}extremely laborious with two magnets, seem likely to present insurmountable practical difficulties for a considerably greater number. The best course will be to abandon this path altogether, which inevitably brings to mind. the attempts to explain planetary motions by piling up more and more epicycles.

In the present treatise I will develop the general theory of terrestrial magnetism independently of all particular hypotheses about the distribution of the magnetic fluid within the Earth, and I will also report the results I obtained from the first application of the method. Imperfect as these results must be, they can still give some idea of what may be hoped for in the future, once reliable and complete observations from every part of the Earth can serve as the basis for refining them more closely and repeatedly.

